\section*{DISCUSSION}

The accumulation of publicly deposited microbiome amplicon sequencing data presents
an enormous opportunity for integrative mega-analysis, yet realising this potential
has been constrained by a set of recurring, interrelated barriers: the labour
intensity of cross-study metadata harmonisation, the systematic exclusion of
CNCB/GSA-archived data from existing retrieval tools, the inability of most
workflows to reconcile datasets generated on disparate sequencing platforms or
targeting different hypervariable regions, and the absence of a single pipeline
that unifies all of these steps end-to-end.
Meta2Data was designed to address each of these barriers within a single,
largely automated workflow requiring only one round of human review throughout the
entire process.

\paragraph{Addressing the metadata integration barrier.}
Metadata harmonisation remains one of the most time-consuming and error-prone
steps in large-scale comparative microbiome studies, primarily because
experimental-design heterogeneity across laboratories produces inconsistent field
definitions and annotation standards even within databases that enforce submission
guidelines~\citep{cernava2022metadata, kim2025fair, bustos2026nonstandardized}.
The MetaDL module substantially reduces this burden by merging SRA RunInfo and
BioSample attribute tables on shared identifiers in a single operation, automatically
appending BioProject descriptions as a dedicated column, and enforcing a uniform
CamelCase column-naming convention via a configurable synonym dictionary.
The companion \texttt{column\_description.tsv} file---reporting per-column fill
rate, dataset coverage, and the five most frequent values for every non-core
field---allows users to rapidly characterise the annotation landscape of an assembled
corpus without inspecting the full merged table.
Primer information, which is frequently absent or inconsistently recorded in
public submissions~\citep{kim2025fair}, is inferred automatically by the downstream
AmpliconPIP module rather than sourced from metadata records, insulating the workflow
from an inherently unreliable data field.
Taken together, these features reduce the single mandatory human-review step to an
exercise in filtering and column selection rather than field-by-field curation,
substantially compressing the time required to prepare a metadata table that is ready
for amplicon processing.

\paragraph{Bridging the CNCB/GSA database coverage gap.}
Perhaps the most consequential limitation of existing re-analysis workflows is
their exclusive targeting of INSDC member databases.
The Genome Sequence Archive (GSA) and GSA-Human portals of the China National Center
for Bioinformation had archived more than 80~PB of data as of August 2025~\citep{cncb2026}, and the
fraction of this archive that represents Asian microbiome cohorts---spanning gut,
oral, and environmental niches that are systematically underrepresented in
NCBI~SRA---makes its omission from integrative analyses a meaningful source of
geographical and biological bias.
Meta2Data is, to our knowledge, the first general-purpose re-analysis pipeline to
support metadata retrieval and sequence download from both NCBI and CNCB/GSA
within a unified interface, accepting BioProject identifiers from all major
international prefixes (PRJNA, PRJEB, PRJDB, and PRJCA) interchangeably and without
any user-side format conversion.
This capability directly closes the coverage gap that affects established tools
including NCBImeta~\citep{eaton2020ncbimeta}, Kingfisher~\citep{woodcroft2024kingfisher},
MADAME~\citep{fumagalli2023madame}, and ffq~\citep{galvez2023ffq}, ensuring that
cross-study analyses can incorporate the full breadth of globally deposited
microbiome data.

\paragraph{Automated cross-platform and cross-region amplicon processing.}
The AmpliconPIP module addresses the platform heterogeneity barrier through two
complementary mechanisms.
First, all sequencing platform assignments are resolved in a single batched Entrez
API query executed before parallel processing begins, preventing the rate-limit
exhaustion that arises when each dataset issues independent platform queries at
runtime.
This cached assignment then routes each dataset to the appropriate denoising
branch---DADA2 paired-end denoising for Illumina (with automatic fallback to
single-end processing when paired-end retention falls below 50\%), DADA2 denoise-ccs
for PacBio~SMRT long-amplicon data, DADA2 denoise-pyro for Ion~Torrent, and
VSEARCH-based OTU clustering for legacy 454 datasets---without requiring any
user-level platform specification.
Second, the positional base-frequency matrix algorithm for primer identification
removes the requirement for users to supply primer sequences, which has been
identified as a persistent impediment to scalable automated re-analysis~\citep{regueira2023critical}.
By classifying each of the first 60 read positions as conserved, degenerate, or
variable and resolving primer termini from continuity of conserved/degenerate blocks,
the algorithm identifies trimming sites with IUPAC-compatible matching across all 16S
hypervariable regions from V1 to V9 and handles mixed-orientation libraries through
bimodal base-distribution analysis.
This design means that a single AmpliconPIP invocation can process datasets sequenced
on different platforms and targeting different variable regions in the same run,
which is a prerequisite for the kind of broad phylogenetic and biogeographical
comparative analyses that motivate mega-analysis in the first place~\citep{xie2024global,cheng2025host}.

\paragraph{Comparison with existing tools.}
Table~\ref{tab:comparison} contextualises Meta2Data relative to the closest
existing tools.
NCBImeta~\citep{eaton2020ncbimeta} provides efficient structured-metadata retrieval
from NCBI databases but performs no sequence download or bioinformatic processing.
Kingfisher~\citep{woodcroft2024kingfisher} offers robust multi-strategy SRA data
procurement but is limited to INSDC databases and produces only raw FASTQ files.
ffq~\citep{galvez2023ffq} and MADAME~\citep{fumagalli2023madame} retrieve metadata
with user-friendly interfaces yet similarly do not support CNCB/GSA access or
downstream amplicon processing.
q2-fondue~\citep{ziemski2022q2fondue}, the closest conceptual predecessor to
Meta2Data, integrates metadata retrieval and SRA data import within the QIIME2
framework; however, it does not support CNCB/GSA data, does not include automatic
primer detection or cross-platform denoising, and requires a QIIME2 installation
even for metadata-only operations.
Meta2Data integrates all of these capabilities---cross-database metadata retrieval
and standardisation, sequence download, platform-aware quality control, automatic
primer detection and trimming, multi-platform denoising, taxonomy classification
against three reference databases, and SEPP-based phylogenetic fragment
insertion~\citep{mirarab2012sepp} into the GreenGenes2 reference
tree~\citep{mcdonald2024greengenes2}---into a single modular pipeline that
outputs QIIME2-compatible feature tables and taxonomy artefacts ready for
downstream diversity analyses without any intermediate format conversion.
Importantly, MetaDL is designed to operate on a bare Python~3 environment without
requiring QIIME2, keeping the metadata curation step lightweight and decoupled from
the amplicon-processing environment.

\paragraph{Limitations.}
Several limitations of the current implementation should be acknowledged.
Meta2Data currently supports Linux operating systems only, restricting direct
usability for Windows or macOS users.
The ggCOMBO module is resource-intensive: the SEPP phylogenetic placement step
requires approximately 70--100~GB of RAM and 20~CPU cores, which in practice
necessitates access to high-performance computing infrastructure for datasets of
non-trivial scale.
The full workflow---from keyword search to final QIIME2 artefacts---can require
one to two days, dominated by sequence download and phylogenetic tree construction.
Oxford Nanopore Technology (ONT) long reads are not currently supported and are
silently skipped.
The automatic primer detection algorithm, while robust across standard 16S library
preparations from V1 to V9, may encounter difficulties with highly non-standard
protocols or with datasets in which primer sequences were removed upstream of
database deposition.
The pipeline is currently restricted to amplicon sequencing data; whole-genome
shotgun (WGS) metagenomic datasets are outside its present scope.
Finally, notwithstanding the substantial automation provided by MetaDL, one round of
human review of the merged metadata table remains necessary to exclude off-target
studies and verify sample-level annotations before amplicon processing, as no
automated system can yet reliably resolve the semantic diversity of free-text
experimental descriptions across heterogeneous BioProject submissions~\citep{bustos2026nonstandardized,kim2025fair}.

\paragraph{Future directions.}
Meta2Data is under active development as an open-source project hosted on GitHub.
A fourth module, ShortreadsPIP, is currently in active development and will extend
the pipeline's scope to whole-genome shotgun metagenomic data, encompassing
read-level quality control, host decontamination, taxonomic profiling, and
functional annotation, thereby enabling unified automated re-analysis of both
amplicon and metagenomic datasets from the same public repositories.
Support for Oxford Nanopore Technology reads is also planned as long-read
sequencing becomes increasingly prevalent in microbiome research.
We anticipate that the BioProject description column automatically populated by
MetaDL---which provides a concise natural-language summary of each study's
purpose, target organism, and experimental design---will serve as a productive
interface for large-language-model-assisted metadata screening, potentially
further reducing the sole remaining manual step and moving large-scale microbiome
re-analysis closer to a fully automated paradigm.
